Classical aerodynamics will study the effects of vortical motion in a fluid otherwise bereft thereof.
A thorough exposition of that matter shall be handled in a later section, though the discussion of the influence vortices on fluids in general shall be accounted for promptly.
We shall in this section examine how the \textsc{Biot}--\textsc{Savart} equation arises as a solution to the problem of a completely vortical fluid.

In order to study the effects of vortices on the velocity field, we define the vorticity $\vorticityfrak = \star\dee\ufrak$, with the relation $\vorticityfrak^{\sharp} = \vorticity$.
We consider the vorticity equation, obtained by taking the curl of the \textsc{Navier}--\textsc{Stokes} equation \subeqref{\ref{eq:Navier-Stokes}}[a].
\begin{subequations}\label{eq:vorticity-equation}
         \begin{equation}\tag{\theequation{a, b}}
           \varrho\deet\vorticityfrak + \varrho\lieder{\uvec}\vorticityfrak = \mu\Delta\vorticityfrak, \qquad \codee\vorticityfrak = 0.
         \end{equation}
\end{subequations}
Confer with \textsc{Batchelor},%
\footnote{%
        \cite{batchelor2007introduction} \textsc{Batchelor}, \S5.2, pp.267--273
}
\textsc{Marsden} \& \textsc{Weinstein},%
\footnote{%
        \cite{marsden1983coadjoint} \textsc{Marsden} \& \textsc{Weinstein} (1983), eq.V
}
or \textsc{Modin} \emph{et al}.%
\footnote{%
        \cite{modin2010geodesics} \textsc{Modin} \emph{et al.} (2010)
}
The question at hand is whether the vorticity is useful in the representation of the fluid flow we are interested in modeling.
For any structure interacting with a fluid inducing a perturbed velocity, this velocity ought to not propagate indefinitely.
In other words, it should be evanescent, a condition of regularity.
Such vector fields may be described by decomposition into constituent \textsc{Helmholtz} components, wherein
\begin{subequations}\label{eq:Helmholtz_decomposition}
        \begin{equation}\tag{\theequation{a, b}}
            \uvec = \nabla \Phi + \nabla \times \mathbf{\Psi}, \qquad \mathbf{\Psi} = \psi \nabla A.
        \end{equation}
\end{subequations}
In terms of the vernacular of differential forms, the so-called scalar potential $\Phi$ is a zero-form, and the vector potential $\mathbf{\Psi}$ is a sharped one-form.
Indeed, as the curl operator is simply $\star\dee$, any choice of bivector $\vfrak$ such that $\star\vfrak = \Psi_i\,\dee x^i$ will be approporiate, revealing the familiar \textsc{Hodge} decomposition of theorem (\ref{thm:Hodge_decomposition}).
Such a choice of representation of the vector potential is called a gauge, and the gauge chosen in equation \subeqref{\ref{eq:Helmholtz_decomposition}}[b] is that of \textsc{Gumerov} \& \textsc{Duraiswami},%
\footnote{%
        \cite{gumerov2006fast} \textsc{Gumerov} \& \textsc{Duraiswami} (2006)
}
where we impose the condition that $\Div\mathbf{\Psi} = 0$.
In other words, the bivector $\vfrak$ is exact.
It ought to be noted that the condition of evanescence is not one which applies to the \textsc{Helmholtz} decomposition, but rather that of the convergence of integrals which will be discussed promptly.

Should the fluid happen to interact with a structure, it is only near the surface the effects of viscosity are apparent.
A formulation, then, wherein the vorticity is nonzero only on the boundary of the object, we may be able to proceed with a potential flow formulation in the rest of the fluid domain.
Nevertheless, the laplacian term in the \textsc{Navier}--\textsc{Stokes} equation may be expanded in terms of equation \eqref{eq:Hodge_laplacian} so as to reveal that
\begin{equation}\label{eq:Laplace_vector_definition}
\mu\Delta \uvec = \mu\nabla(\nabla \cdot \uvec) - \mu \nabla \times \vorticity.
\end{equation}
This immediately yields the biharmonic equation for the potential function $\Phi$ when the vorticity is zero, which has been discussed by \textsc{Gumerov} and \textsc{Duraiswami} in terms of implementation through fast multipole methods.\footnote{\emph{ibid}.}
For we know the divergence of the velocity field to be zero, so that the Laplacian is equal to the curl of the vorticity.
The connexion is immediately apparent between that of viscous effects encoded in the Laplacian term, and the vorticity of the fluid.
Thus our motivation is plainly laid out before us---we wish to imitate the viscous effects that the fluid exhibits through a potential formulation.
Employing the above \textsc{Helmholtz} decomposition of equation \eqref{eq:Helmholtz_decomposition}, we find that by taking the laplacian thereof and substituting in the value of the vorticity of equation \eqref{eq:Laplace_vector_definition},
\[
  -\nabla \times \vorticity = \nabla \times \Delta \mathbf{\Psi}.
\]
We are free to equate the two arguments of the curl so that the vorticity equals the negative laplacian of the vector field $\mathbf{\Psi}$, or in terms of differential forms, $\Delta\star\vfrak = -\vorticityfrak$.
This is a \textsc{Poisson} equation, which is solved by the convolution $\star\vfrak = -\mathfrak{w} \ast \green$,%
\footnote{%
        \cite{taylor2023partial} \textsc{Taylor}, \S3.4
}
where $\green$ is the familiar \textsc{Green} function, expressed in local coördinates as
\[
        \Psi_i(\xvec) = -\int_{\Omega(\bzhe)} w_i \green \,\dee V, \qquad \green(\xvec, \bzhe) = \frac{1}{4\pi \absl{\xvec - \bzhe}}.
\]
Since it is the velocity $\ufrak = \Div\vfrak$ we wish to find, we may simply calculate it from the above expression by taking the curl.
Note that integration is done with respect to the variable $\bzhe$, whilst the exterior derivative will be taken with respect to $\xvec$.
To that end, we have $\partial_i \vorticityfrak(\bzhe)\dee x^{i} = 0$, so that by the graded product rule, we only differentiate the \textsc{Green} function.
We then find the velocity vector field in terms of the vorticity, such that
\begin{equation}
        \ufrak(\xvec) = -\int_{\Omega} \star\dee(\vorticityfrak \green) = -\int_{\Omega} \star\big( \mathfrak{w}\wedge\dee\green \big).
\end{equation}

\noindent
The above expression seems to hold water,%
\footnote{%
        \cite{tyssvang2025derivation} \textsc{Tyssvang} (2025)
}
as a similar result was obtained by \textsc{Hesin}.%
\footnote{%
        \cite{hesin2012symplectic} \textsc{Hesin} (2012)
}
The proof in the latter is somewhat different and more rigorously involved, as the vorticity is given as the two-form $\star\vorticityfrak$, yielding a componentwise \textsc{Poisson} equation $\Delta u^i = \star(\dee x^{i} \wedge \star \codee \star \vorticityfrak)$.

As was noted above, the condition of evanescence of the velocity is a result of this formulation.%
\footnote{%
        cite
}

We note the relation between the \textsc{Hodge} star of the exterior product of two 1-forms and the cross product of two vectors; $\star(\afrak \wedge \bfrak) = (\avec \times \bvec)$, up to musical isomorphism.
Expressing this in terms of vectors, we have the generalized \textsc{Biot}--\textsc{Savart} equation,
\begin{equation}
        \uvec(\xvec) = -\frac{1}{4\pi}\int_{\Omega} \frac{ \vorticity \times (\xvec - \bzhe) }{ {\absl{ \xvec - \bzhe }}^3 } \,\dee V.
\end{equation}


Introduce circulation, such that the vorticity is given along vortex filaments, $\vorticityfrak = \Gamma\mathfrak{t}$.
Gives compact support, etc. so integrals converge nicely.
Provide definition of circulation.
\[
  \Gamma = \oint_{C} \uvec \cdot \dee\mathbf{s} = \int_{S} \vorticity \cdot \nhat \,\dee S.
\]
We can derive the circulation theorem.
\[
  \Dee_t \Gamma = 0.
\]
We will derive how boundary layers may be approximated by vortex sheets.

The \textsc{Biot}--\textsc{Savart} formulation of the disturbance velocity field at any point by the filamtent $C$ is then given by%
\footnote{%
        \cite{deng2020acceleration} \textsc{Deng} \emph{et al.} (2020), eq.6
}%
\textsuperscript{,}%
\footnote{%
        \cite{schlichting1967aerodynamiki} \textsc{Schlichting} \& \textsc{Truckenbrodt}, p.98, eq.2.161
}
\begin{equation}\label{eq:Biot-Savart}
  \uvec(\bzhe) = \frac{1}{4\pi} \int_{C} \Gamma(\xvec)\frac{\xvec - \bzhe}{r^3} \times \dee\mathbf{s}.
\end{equation}

\noindent
